% sec-02: the five letters, the signing order, and the single approval step.
% Figure 6 (mermaid-type sequence), Table 7, and the actionbox.
\section{The Action Register}

Five letters leave on this day and they are deliberately not interchangeable. One
opens a commercial relationship that has never been opened. Two ask for
introductions and describe no security. One engages counsel on a defined scope.
One prepares an account that would receive a subscription that does not yet
exist.

\begin{apptable}
\begin{tabularx}{\textwidth}{>{\raggedright\arraybackslash}p{0.5cm} >{\raggedright\arraybackslash}p{3.4cm} >{\raggedright\arraybackslash}X}
\headrow \thead{\#} & \thead{Recipient} & \thead{The ask, and what the letter deliberately does not do}\\
\midrule
1 & The agent's developer, external research & A scientific feasibility conversation and the next submission window. It does not ask for drug, funding, or a letter of authorization in this letter \\
2 & Regional angel and syndicate screening committees & A screening conversation. It states no valuation, cap, discount, minimum, or closing date \\
3 & Life science family offices & A read of the public repository, then a call. It carries no attachment at all \\
4 & Securities counsel & An engagement scoped to one instrument, two filings, and one eligibility question. It does not state a conclusion the company has relied on \\
5 & The broker-dealer corporate desk & A subscription account opened and left unfunded. It states plainly that no offering exists \\
\end{tabularx}
\end{apptable}
\vspace{-0.60cm}
\tabcap{Table 7. The five letters of this day and the ask carried in each,\\
beside the thing every one of them is deliberately written not to do.}

\subsection{The order, and the idleness inside it}

Figure 6 draws the signing sequence of a financing across four participants. Its
useful content is not the order, which is obvious, but the idle spans: counsel is
idle for four consecutive spans in the middle, and a purchaser is idle for seven.
That distribution is the argument for engaging counsel on a defined scope rather
than an open retainer. The scope is what counsel does in the spans where counsel
is not idle, and it is priceable in advance.

\begin{appfloat}[!tb]
\begin{appfig}[mermaid-type / sequence with lifelines and activation bars]
% Four lifelines on a 3.30cm pitch; nine messages on a 0.62cm pitch.  \seqrow
% draws one message row, so a message can be moved by changing one number.
\node[mmactor] (ceo) at (0,0)    {Chief executive};
\node[mmactor] (cns) at (3.30,0) {Securities counsel};
\node[mmactor] (inv) at (6.60,0) {Purchaser};
\node[mmactor] (brk) at (9.90,0) {Broker-dealer};
\foreach \n in {ceo,cns,inv,brk}{\draw[mmlife] (\n.south) -- ++(0,-6.05);}
\seqrow{-0.85}{0}{3.30}{mmmsg}{Engagement scope, four items}
\seqrow{-1.47}{3.30}{0}{mmret}{SBIR ownership answer, in writing}
\seqrow{-2.09}{0}{3.30}{mmmsg}{Instrument selected}
\seqrow{-2.71}{3.30}{0}{mmret}{Documents drafted, filings prepared}
\seqrow{-3.33}{0}{9.90}{mmmsg}{Open the subscription account}
\seqrow{-3.95}{9.90}{0}{mmret}{Account open, wire instructions issued}
\seqrow{-4.57}{0}{6.60}{mmmsg}{Documents, after a prior relationship exists}
\seqrow{-5.19}{6.60}{9.90}{mmmsg}{Subscription funds}
\seqrow{-5.81}{9.90}{0}{mmret}{Funds received, first sale dated}
% Activation bars on the participant that owns each span.
\node[mmact,minimum height=1.9cm] at (3.30,-2.10) {};
\node[mmact,minimum height=1.0cm] at (9.90,-3.65) {};
\node[mmact,minimum height=1.0cm] at (6.60,-4.90) {};
\node[mmact,minimum height=5.3cm] at (0,-3.35) {};
\pnote{10.60}{-2.10}{Counsel idle after\\the written answer}
\pnote{10.60}{-4.60}{Purchaser appears\\in two spans only}
\end{appfig}
\vspace{-0.60cm}
\figcaption{Figure 6. The signing order of a private financing across four participants,\\
drawn so that both the sequence and the idle spans between it are visible.}
\end{appfloat}

\subsection{Why the developer letter is on this day and not the last one}

The developer's public external research portal has indicated that submissions
were closed \cite{revmedexternal2026}. The approval changes this company's
posture, not that portal's calendar, so the letter asks for the next submission
window and a scientific conversation and asks for nothing a closed window would
have to refuse.

It sits after the federal re-contacts of day 1 because a developer reasonably
asks who is funding the work, and day 1 is the answer to that question. It sits
before the site approaches of day 3 because a site asks the same question in a
different form and then asks who the investigator is.

\subsection{What is deliberately not sent today}

No approach is made to a clinical site on this day, and none is made to a
foundation. A site asks who is funding the work and then asks who the
investigator will be; the first question is answered by day 1 and the second by
the firewall in \S4, and neither answer is improved by being given on the same
day the capital question is still open. Foundations run on published cycles
rather than on correspondence, so their approach belongs on a day built around
cycles rather than around instruments. Both are day 3.

\subsection{The single approval step}

\actionbox{One decision is open on this day}{Approve one instrument from \S3, one
raise size, and the first approach to the agent's developer in
\appfile{funding/auto-fund/03Sep26/emails}. The comparison is finished, the
letters are written, and the two filings are prepared and deliberately unfiled.
Letters 2 and 3 must not be sent to more than one organization on a single
recipient line, and no letter may state an offering term.}
